\section{Sharp partial identification of ranking metrics}
\label{sec:ranking}

Existing sensitivity analysis handles functionals that are
\emph{linear-fractional} in the outcome regression --- sensitivity and
specificity at a threshold, predictive values, treatment effects
\citep{zhao2019sensitivity, dorn2023sharp, yadlowsky2018bounds} --- for which
sharp bounds follow from a one-dimensional search. AUROC, the metric that
governs model selection in our applications, is not of that form. Write
$\psi(u,v) = \one\{u>v\} + \tfrac12\one\{u=v\}$ and $\mu = P_X$. For a score
$f$, regression $p$ and prevalence $\pi = \E_\mu[p]$,
\begin{equation}
  \AUC(p)
  \;=\;
  \frac{\E_{\mu\otimes\mu}\big[\psi\big(f(X),f(X')\big)\,p(X)\big(1-p(X')\big)\big]}
       {\pi\,(1-\pi)} .
  \label{eq:auc-def}
\end{equation}
The numerator is bilinear in $p$ and the denominator quadratic, so $\AUC$ is
neither convex nor concave over the box of \cref{lem:box}: AUROC couples every
unit to every other, and the pointwise decoupling behind threshold metrics
fails.

The obstruction dissolves through an identity. The \emph{mid-rank function} of
$f$ under $\mu$,
\begin{equation}
  \begin{aligned}
  R(x) &= \E_{X'\sim\mu}\big[\psi\big(f(x), f(X')\big)\big] \\
       &= \mu\big(f(X') < f(x)\big) + \tfrac12\,\mu\big(f(X')=f(x)\big),
  \end{aligned}
  \label{eq:midrank}
\end{equation}
is \emph{identified}: it depends only on the score and on $P_X$.

\begin{theorem}[Rank identity]
\label{thm:rank-identity}
For every measurable $p:\X\to[0,1]$ with $\pi=\E_\mu[p]\in(0,1)$,
\begin{equation}
  \boxed{\;
  \AUC(p)
  \;=\;
  \frac{\E_\mu\!\big[\,p(X)\,R(X)\,\big] \;-\; \pi^{2}/2}{\pi\,(1-\pi)} \; }
  \label{eq:rank-identity}
\end{equation}
In particular, \emph{at fixed prevalence $\pi$, AUROC is an affine functional of
the outcome regression.}
\end{theorem}

The proof (\cref{app:ranking}) is two lines: because
$\psi(u,v)+\psi(v,u)=1$ identically, the kernel operator
$(\mathcal{K}h)(x) = \E_\mu[\psi(f(x),f(X'))\,h(X')]$ satisfies
$\mathcal{K}+\mathcal{K}^{*} = \one\otimes\one$, so the quadratic term
$\langle p, \mathcal{K}p\rangle$ equals $\pi^2/2$ for \emph{every} $p$ and
collapses into a function of prevalence alone. \Cref{eq:rank-identity} is a
population form of the Mann--Whitney representation
\citep{mann1947test, hanley1982meaning}; we have not seen it used to linearise
AUROC in the outcome regression.

\subsection{The sharp interval, exactly, in \texorpdfstring{$O(n\log n)$}{O(n log n)}}

\Cref{thm:rank-identity} makes the sharp interval a two-stage problem:
optimise the linear functional $\langle p,R\rangle$ over the box at fixed
prevalence, then scan the prevalence.

\begin{theorem}[Sharp AUROC interval]
\label{thm:sharp-auc}
Let $\I_\Gamma$ be the box of \cref{lem:box}, let
$\pi_{\!-}=\E_\mu[\plo]$, $\pi_{+}=\E_\mu[\pup]$, and for
$\pi\in[\pi_-,\pi_+]$ define
\[
  V^{\pm}(\pi) \;=\; \genfrac{}{}{0pt}{}{\max}{\min}
  \Big\{\, \E_\mu[p\,R] \;:\; p\in\I_\Gamma,\ \E_\mu[p]=\pi \,\Big\}.
\]
Then
\begin{equation}
  \Big[\;\inf_{p\in\I_\Gamma}\AUC(p),\ \sup_{p\in\I_\Gamma}\AUC(p)\;\Big]
  \;=\;
  \Big[\;\min_{\pi\in[\pi_-,\pi_+]}\frac{V^{-}(\pi)-\pi^2/2}{\pi(1-\pi)},\ \
          \max_{\pi\in[\pi_-,\pi_+]}\frac{V^{+}(\pi)-\pi^2/2}{\pi(1-\pi)}\;\Big],
  \label{eq:sharp-auc}
\end{equation}
and both endpoints are \emph{attained} by members of $\I_\Gamma$. Moreover:
\begin{enumerate}
\item[(i)] $V^{+}$ is piecewise linear and concave, $V^{-}$ piecewise linear and
  convex, each with at most $n$ breakpoints; both are computed by a continuous
  knapsack --- fill $p$ up from $\plo$ toward $\pup$ in decreasing (resp.\
  increasing) order of $R(x)$ until the budget $\pi$ is exhausted.
\item[(ii)] On each linear piece $V^\pm(\pi)=\alpha+\beta\pi$ the objective is a
  ratio of quadratics whose stationary points solve
  $(\beta-\tfrac12)\pi^{2} + 2\alpha\pi - \alpha = 0$; the cubic terms cancel
  identically.
\item[(iii)] Consequently \eqref{eq:sharp-auc} is computed \emph{exactly} in
  $O(n\log n)$ time and $O(n)$ space.
\end{enumerate}
\end{theorem}

\begin{proposition}[A unified class]
\label{prop:unified}
\Cref{thm:sharp-auc} holds verbatim for every functional of the form
\[
  M(p) \;=\;
  \frac{\E_\mu[p\,\phi] + c_0 + c_1\pi + c_2\pi^{2}}
       {d_0 + d_1\pi + d_2\pi^{2}},
\]
with $\phi$ identified, the stationary condition becoming
$A\pi^2+B\pi+C=0$ with $A=c_2d_1-c_1'd_2$, $B=2(c_2d_0-c_0'd_2)$,
$C=c_1'd_0-c_0'd_1$ where $c_0'=\alpha+c_0$, $c_1'=\beta+c_1$. Taking
$\phi=R,\,c_2=-\tfrac12,\,d_1=1,\,d_2=-1$ gives AUROC; $\phi=\one\{f>\tau\}$,
$d_1=1$ gives sensitivity at threshold $\tau$; $\phi=-\one\{f\le\tau\}$,
$c_0=\mu(f\le\tau)$, $d_0=1,d_1=-1$ gives specificity; $\phi=\one\{f>\tau\}$,
$d_0=\mu(f>\tau)$ gives PPV. The known fixed-threshold results are thus special
cases, and AUROC is the new member.
\end{proposition}

\subsection{Why the obvious recipe is wrong}

The natural move is to evaluate the metric at the corners of $[\plo,\pup]$ and
report the range. For threshold metrics that is correct; for AUROC it is
\emph{not valid}:

\begin{proposition}[Corner evaluation under-covers]
\label{prop:corner}
Let $\mathcal{C} = [\min,\max]$ of $\{\AUC(\plo),\AUC(\tilde p),\AUC(\pup)\}$.
Then $\mathcal{C}$ is contained in the sharp interval \eqref{eq:sharp-auc}, and
the containment is strict whenever $R$ is non-constant on
$\{x:\delta(x)>0\}$. Consequently there exist laws in $\I_\Gamma$ whose AUROC
lies outside $\mathcal{C}$, and $\mathcal{C}$ fails to cover the truth.
\end{proposition}

The cause is exactly the coupling \cref{thm:rank-identity} exposes: the
extremal member of $\I_\Gamma$ sets $p=\pup$ where the score ranks units
\emph{high} and $p=\plo$ where it ranks them low, with a single fractional unit
in between --- a mixed vertex no corner evaluation visits.
Empirically (\cref{sec:experiments}) the corner interval covers the true
deployment AUROC in $10\%$ of benchmark cells versus $100\%$ for the sharp
interval, while being $14\times$ narrower --- narrow because it is wrong.

A valid but loose alternative is to bound numerator and denominator of
\eqref{eq:rank-identity} separately. That interval always contains the sharp one
and in our benchmarks is $3.2\times$ wider, which quantifies what sharpness buys.
